% ID pendamping stabil: methods-of-algebra-volume-2-id/Errata-Al-jabr-2-id | Sumber: authority/upstream/AlJabr-2-9a5803ff77dd3257484cb177f851a73770a59dd3/Errata-Al-jabr-2.tex | Cakupan: berkas lengkap, baris 1--208 | Lisensi: CC BY 4.0.
%!TEX TS-program = xelatex
%!TEX encoding = UTF-8

% Sumber LaTeX untuk ralat buku “Metode dalam Aljabar” (Jilid 2), diterjemahkan ke bahasa Indonesia dari bahasa Tionghoa
% Copyright 2026  李文威 (Wen-Wei Li).
% Permission is granted to copy, distribute and/or modify this
% document under the terms of the Creative Commons
% Attribution 4.0 International (CC BY 4.0)
% http://creativecommons.org/licenses/by/4.0/

% Ralat Metode dalam Aljabar, Jilid 2 / Wen-Wei Li
% Menggunakan kelas dokumen khusus AJerrata.cls. Memuat xeCJK secara otomatis.

\documentclass{AJerrata}

\usepackage{unicode-math}

\usepackage[unicode, colorlinks, psdextra, bookmarksnumbered,
	pdfpagelabels=true,
	pdfauthor={李文威 (Wen-Wei Li)},
	pdftitle={Ralat Metode dalam Aljabar, Jilid 2},
	pdfkeywords={}
]{hyperref}

\setmainfont[
	BoldFont={texgyretermes-bold.otf},
	ItalicFont={texgyretermes-italic.otf},
	BoldItalicFont={texgyretermes-bolditalic.otf},
	PunctuationSpace=2
]{texgyretermes-regular.otf}

\setsansfont[
	BoldFont=FiraSans-Bold.otf,
	ItalicFont=FiraSans-Italic.otf
]{FiraSans-Regular.otf}

\setCJKmainfont[
	BoldFont=Noto Serif CJK SC Bold
]{Noto Serif CJK SC}

\setCJKsansfont[
	BoldFont=Noto Sans CJK SC Bold
]{Noto Sans CJK SC}

\setCJKfamilyfont{emfont}[
	BoldFont=FandolHei-Regular.otf
]{FandolHei-Regular.otf}	% Fonta untuk penekanan

\renewcommand{\em}{\bfseries\CJKfamily{emfont}} % Penekanan

\setmathfont[
	Extension = .otf,
	math-style= TeX,
]{texgyretermes-math}

\usepackage{mathrsfs}
\usepackage{stmaryrd} \SetSymbolFont{stmry}{bold}{U}{stmry}{m}{n}	% Menghindari peringatan (stmryd tidak memiliki huruf tebal)
% \usepackage{array}
% \usepackage{tikz-cd}  % Menggunakan TikZ untuk menggambar
\usetikzlibrary{positioning, patterns, calc, matrix, shapes.arrows, shapes.symbols}

\usepackage{myarrows}				% Menggunakan panah dapat-renggang khusus
\usepackage{mycommand}				% Memuat perintah kebiasaan khusus


\title{\bfseries Ralat Metode dalam Aljabar (Jilid 2) \\ Rentang: sejak terbit resmi pada September 2024 hingga kini }
\author{李文威 (Wen-Wei Li)}
\date{\today}

\begin{document}
	\maketitle

	\begin{Errata}
		\item[Definisi 1.5.2, butir kedua (memantulkan limit) dan butir ketiga (menciptakan limit)]
		Hapus ``(jika ada)'' dari butir kedua.
		
		Ganti ``Jika ...'' pada butir ketiga beserta daftar sesudahnya dengan kalimat berikut: ``... jika $\varinjlim F\alpha$ ada, maka $\varinjlim \alpha$ ada, dan dalam hal ini $F$ mempertahankan sekaligus memantulkan $\varinjlim$ ini.''
		\Thx{Terima kasih kepada 王昱翔 (Wang Yuxiang) atas koreksinya}
		
		\item[Proposisi 1.5.7, paragraf kedua terakhir dalam pembuktian]
		\Orig ... diangkat menjadi $L \to X$.
		\Corr ... diangkat menjadi $X \to L$.
		
		\item[Pembahasan tepat di atas Catatan 1.7.3]
		\Orig (iv)
		\Corr (ii)
		
		\item[Lema 1.9.8, subskrip $\varinjlim$ pada ruas kanan rumus yang ditampilkan dalam pernyataan, serta paragraf terakhir pembuktian]
		Ganti $S_{/Y}$ dengan $S_{Y/}$ pada kedua tempat.
		
		\item[Korolari 1.11.14, baris pertama pembuktian]
		\Orig $\alpha: I \to \mathcal{C}$
		\Corr $\alpha: I \to \mathcal{D}$
		
		\item[Latihan 8 Bab 1]
		Ganti subskrip $L$ pada $\mathrm{Lan}$ dan $\mathrm{Ran}$ dengan $K$ (dua tempat).
		
		\item[Petunjuk Latihan 11 Bab 1] Dalam paragraf terakhir ``Untuk (iv) $\implies$ (ii) ...'', ganti ``adjoin kiri'' dengan ``adjoin kanan'' (dua tempat).
		
		\item[Petunjuk Latihan 12 Bab 1] Dalam diagram kedua pada petunjuk, ganti $t^\dagger$ dan $b^\dagger$ masing-masing dengan $t$ dan $b$.
		
		\item[Pernyataan Teorema 2.4.6]
		Dalam rumus yang ditampilkan pada baris setelah ``terdapat isomorfisme standar antarinterval ...'', ganti $(u^{\flat} \wedge v) \vee (v \wedge v^{\flat})$ dalam interval dengan $(u^{\flat} \wedge v) \vee (u \wedge v^{\flat})$.
		
		Dalam kalimat terakhir ``semuanya direalisasikan oleh pemetaan dalam Proposisi 2.4.5 ...'', ganti $(u \wedge v)$ dengan $(u \wedge v^{\flat})$.
		
		\item[Konvensi 2.6.3, baris kedua]
		\Orig supremum (atau infimum)
		\Corr infimum (atau supremum)
		\Thx{Terima kasih kepada 黄行知 (Huang Xingzhi) atas koreksinya}
		
		\item[Teorema 3.2.6, butir kedua pernyataan]
		\Orig homotopi $g$ dari $f$ ke $g$
		\Corr homotopi $h$ dari $f$ ke $g$
		\Thx{Terima kasih kepada 李卓睿 (Li Zhuorui) atas koreksinya}
		
		\item[Paragraf pertama \S 3.12]
		\Orig ... funktor turunan kiri (atau funktor turunan kanan);
		\Corr ... funktor turunan kanan (atau funktor turunan kiri);
		\Thx{Terima kasih kepada 黄行知 (Huang Xingzhi) atas koreksinya}
		
		\item[Korolari 3.12.7, rumus yang ditampilkan pada baris kedua terakhir pembuktian]
		Ganti suku terakhir $\mathrm{R}^1 F(Z)$ dengan $\mathrm{R}^1 F(X)$.
		\Thx{Terima kasih kepada 黄行知 (Huang Xingzhi) atas koreksinya}
		
		\item[Konvensi 3.12.8]
		\Orig funktor turunan kiri berderajat tinggi (atau funktor turunan kanan)
		\Corr funktor turunan kanan berderajat tinggi (atau funktor turunan kiri)
		\Thx{Terima kasih kepada 黄行知 (Huang Xingzhi) atas koreksinya}
		
		\item[Pembuktian Proposisi 3.13.13]
		Dalam paragraf yang dimulai dengan ``Sekarang masuk ke pokok pembahasan ...'', ganti $\psi^{-1}(c)$ di bagian akhir dengan $(\varprojlim \psi)^{-1}(c)$.
		
		\item[Paragraf sebelum Catatan 3.14.8]
		\Orig $\cdots \to Q_1 \to Q_0 \to X \to 0$
		\Corr $\cdots \to Q_1 \to Q_0 \to Y \to 0$
		
		\item[Paragraf keempat terakhir dalam \S 3.14]
		\Orig Sebagai akibatnya, ... $\Hm^p(C) \otimes \Hm^q(D)$, dan dari sana ke $\Hm^n(C \otimes D)$ ...
		\Corr Sebagai akibatnya, ... $\Hm_p(C) \otimes \Hm_q(D)$, dan dari sana ke $\Hm_n(C \otimes D)$ ...
		\Thx{Terima kasih kepada 黄行知 (Huang Xingzhi) atas koreksinya}
		
		\item[Latihan 24 Bab 3]
		Dalam kalimat pertama, ganti ``$\varinjlim$ kecil'' dengan ``$\varinjlim$ kecil terfilter''.
		
		Hapus seluruh kalimat ``Demikian pula, dalam kategori Abel ...'' beserta rumus di bawahnya, lalu ganti dengan ``Bahas apakah terdapat kesimpulan serupa untuk funktor $\Ext$.''
		
		\item[Catatan 4.1.11, baris pertama]
		\Orig funktor kohomologis $F$
		\Corr funktor kohomologis $H$
		\Thx{Terima kasih kepada 张润程 (Zhang Runcheng) atas koreksinya}
		
		\item[Definisi 4.5.11, baris ketiga]
		\Orig ... $X$ isomorfik dengan $Y$ ...
		\Corr ... $X$ melalui $Y$ ...
		\Thx{Terima kasih kepada 郑维喆 (Zheng Weizhe) atas koreksinya}
		
		\item[Teorema 4.5.13, paragraf kedua terakhir pembuktian]
		Paragraf semula ``Untuk bagian kanan diagram ... sementara keadaan pada bagian kiri tetap tidak berubah (silakan verifikasi).'' harus diubah menjadi bentuk berikut:
		
		``Untuk bagian kanan diagram, berdasarkan $a_i^{-n} = \identity_Y$, morfisme $W \to Z_i$ dapat dimodifikasi dengan suatu homotopi yang sesuai agar bagian kanan komutatif di dalam $\cate{C}(\mathcal{A})$. Kemudian, ulangi konstruksi pushout dalam pembuktian surjektivitas untuk mereduksi ke kasus $W^{-n} = Y$ dan $b^{-n} = \identity_Y$, sementara keadaan pada bagian kiri tetap tidak berubah (silakan periksa).''
		\Thx{Terima kasih kepada 黄良伟 (Huang Liangwei) atas koreksinya}
		
		\item[Definisi 5.1.1]
		Ganti $\mathrm{F}^{p+1}$ pada butir pertama dengan $\mathrm{F}^{p+1} X$, dan ganti kategori $\mathrm{F}_{\bullet}(\mathcal{A})$ di akhir paragraf kedua setelah definisi dengan kategori $\mathrm{Fil}_{\bullet}(\mathcal{A})$.
		
		\item[Korolari 5.5.6, baris kedua terakhir pernyataan]
		\Orig dan $\mathrm{F}^n X = X$
		\Corr dan $\mathrm{F}_n X = X$
		\Thx{Terima kasih kepada 黄行知 (Huang Xingzhi) atas koreksinya}
		
		\item[Latihan 8 (i) Bab 5]
		\Orig berdegenerasi pada halaman $E_{\mathrm{I}}^1$.
		\Corr berdegenerasi pada halaman $E_{\mathrm{I}}^2$.
		
		\item[Pembuktian Lema 7.1.8]
		Ganti $\eta_A \otimes_B$ dalam diagram dengan $\eta_A \otimes \eta_B$.
		\Thx{Terima kasih kepada 黄良伟 (Huang Liangwei) atas koreksinya}
		
		\item[Contoh 7.1.11, syarat kedua dalam daftar terakhir]
		\Orig $\otimes: \mathcal{A} \times \mathcal{A} \to \mathcal{A}$ merupakan funktor aditif dalam setiap variabel.
		\Corr $\otimes: \mathcal{A} \times \mathcal{A} \to \mathcal{A}$ merupakan funktor aditif yang mempertahankan koproduk tersebut dalam setiap variabel.
		\Thx{Terima kasih kepada 黄良伟 (Huang Liangwei) atas koreksinya}
		
		\item[Rumus tampilan kedua dalam \S 7.2]
		Ganti $(d_M x) + (-1)^a x (d_N y)$ dengan $(d_M x) \otimes y + (-1)^a x \otimes (d_N y)$.
		\Thx{Terima kasih kepada 黄良伟 (Huang Liangwei) atas koreksinya}
		
		\item[Lema 7.6.12, diagram dalam paragraf kedua pembuktian, di bawah ``... sehingga diperoleh diagram komutatif pada bagian bergaris penuh'']
		Ganti $Tu$ dan $Tv$ di atas dan di bawah panah ganda pada baris kedua masing-masing dengan $u$ dan $v$.
		\Thx{Terima kasih kepada 郝尚诚 (Hao Shangcheng) atas koreksinya}
		
		\item[Rumus pertama dan kedua yang ditampilkan di atas (7.9.3)]
		Ganti $\tilde{a}_{ij}$ dengan $\tilde{a}_{ji}$ (dua tempat).
		
		\item[Baris tepat di atas Catatan 7.9.7]
		Ganti $\tilde{a}_{ij}$ dengan $\tilde{a}_{ji}$.
		
		\item[Definisi--Proposisi 9.1.3, baris pertama pembuktian]
		\Orig teks Tionghoa ``仅伦'' (salah ketik)
		\Corr teks Tionghoa ``仅论'' (``hanya membahas'')
		
		\item[Definisi A.2.11, paragraf terakhir]
		\Orig Ketika $\kappa$ semakin besar, syaratnya menjadi semakin longgar, ...
		\Corr Ketika $\kappa'$ cukup besar relatif terhadap $\kappa$, syarat-syarat yang bersesuaian menjadi lebih longgar daripada syarat yang berlaku untuk $\kappa$; ...
		\Thx{Terima kasih kepada 黄行知 (Huang Xingzhi) atas koreksinya}
	\end{Errata}
\end{document}
